\chapter{Reproduction}
\label{app:repro}

\section{Running the system}
\begin{lstlisting}[language=bash,basicstyle=\ttfamily\scriptsize]
# Tools: Icarus Verilog 12.0; for timing and area also Yosys, OpenSTA (sta)
# and the sky130 liberty file (see backend/pdk/README.md)
cd backend
export PATH="$HOME/.local/bin:$PATH"

# Offline demo (built-in designs, no model calls)
OFFLINE_MODE=1 python3 main.py

# Live backend through the OpenCode app; model keys for Cloudflare
# Workers AI and NVIDIA NIM are read only from backend/.env
env -u OFFLINE_MODE LLM_PROVIDER=opencode-app \
  LLM_MODEL=deepseek-v4.1-flash-free python3 main.py

# Any other OpenAI-compatible endpoint (example: OpenRouter)
env -u LLM_PROVIDER -u OPENCODE_API_KEY -u ZEN_API_KEY -u NVIDIA_API_KEY \
  OPENAI_API_KEY=... OPENAI_BASE_URL=https://openrouter.ai/api/v1 \
  LLM_MODEL=nvidia/nemotron-3-ultra-550b-a55b:free python3 main.py

# Frontend dev server
cd ../frontend && npm run dev -- --host 127.0.0.1
\end{lstlisting}

\section{Verifying the tool}
\label{app:verify}
These commands reproduce Table~\ref{tab:verify}. All except the last make no model
calls.
\begin{lstlisting}[language=bash,basicstyle=\ttfamily\scriptsize]
cd backend
export PATH="$HOME/.local/bin:$PATH"
OFFLINE_MODE=1 python3 -m pytest            # 453 passed
python3 tools/mutation_sweep.py             # 27/27 injected bugs caught
python3 tools/verify_mapping.py             # 19/19 designs, same results at gate level
python3 evaluation/independent/run.py       # 14/14 evaluated designs correct
python3 evaluation/independent/strength.py  # strength of the reference testbenches
python3 tools/preflight.py --live           # one small model call; fallback ladder

cd ../frontend
npx tsc -b && npx vitest run && npm run build   # 222 passed
\end{lstlisting}

The live evaluation of Section~\ref{sec:eval} is re-run with
\code{python3 tools/evaluate_live.py} (it resumes and skips prompts that already
have a record); \code{--summary} only rebuilds \code{evaluation/summary.md} from the
saved records in \code{evaluation/results.jsonl}.

The live harness used for the earlier priority-encoder and seeded-bug runs calls
the same event generator as the web API (\code{pipeline_events}). The resulting RTL,
testbench, simulator transcript, and structured diagnosis are archived under
the directory \path{docs/evidence/} for audit.


\chapter{Complete Prompt Set}
\label{app:prompts}
The system prompts below are extracted from \code{backend/llm_service.py} (only the dashes are changed for typesetting). The same testbench prompt is used for a separate testbench model and for the second-opinion testbench.
\lstset{basicstyle=\ttfamily\tiny,language=,breaklines=true,frame=single}
\section{Specification parser}
\lstinputlisting{appendix/prompt_intent.txt}
\section{RTL generator}
\lstinputlisting{appendix/prompt_rtl.txt}
\section{Independent testbench generator}
\lstinputlisting{appendix/prompt_testbench.txt}
\section{Debugger (diagnosis and repair)}
\lstinputlisting{appendix/prompt_fix.txt}

\chapter{Live-Generated Testbench}
\label{app:tb}
The specification-derived testbench generated live for the priority encoder of Chapter~\ref{sec:results} (archived as \code{docs/evidence/livepriority08_testbench.sv}), followed by the simulator transcript.
\lstinputlisting[style=vlog,basicstyle=\ttfamily\tiny]{evidence/livepriority08_testbench.sv}
\lstinputlisting[language=,basicstyle=\ttfamily\tiny]{evidence/livepriority08_transcript.log}
